\subsection{Irreducible supports with a class VII resolution}

\begin{lemma}[C.2.VII.1; Proved anticanonical support comparison]
\label{st:C.2.VII.1}
Let $Z$ be a smooth compact surface of class VII, not
necessarily minimal, and let $D>0$ be an effective integral
divisor such that $K_Z+D$ is torsion in $\Pic(Z)$.
Write $R=D_{\mathrm{red}}$, and let $r(R)$ be its number
of rational irreducible components. Then
\[
 e(R)=r(R),\qquad r(Z)\le b_2(Z),
\tag{C.VII.1}
\]
where $r(Z)$ counts all rational curves of $Z$.
Every connected component of $R$ has arithmetic genus one.
In particular a connected component containing an elliptic
curve consists of that single elliptic curve with trees of
rational curves attached; it cannot contain a second
elliptic curve or a path joining two distinct points of the
first one.
\end{lemma}
\begin{proof}
We specify the established surface results used in this
comparison. For a minimal class VII surface with positive
second Betti number, existence of a numerically
anticanonical divisor implies existence of a global
spherical shell by
\cite[Theorem~0.2, Proposition~2.29 and Theorem~2.38]{Dloussky2006}.
The reduced support is one rational cycle with attached
trees, one or two rational cycles, or, in the parabolic
Inoue case, a disjoint rational cycle and a smooth elliptic
curve. The intermediate graph and its actual curve
configuration are described in
\cite[\S1, pp.~284--285, and \S3, p.~296]{DOT2003}.
The fact that no connected piece of the anticanonical
support is an isolated tree is
\cite[Corollary~2.25 and Proposition~2.27]{Dloussky2006};
Lemma~2.22 there shows that a curve meeting that support
belongs to it. These statements apply to our divisor with
index one, since its twisting line is torsion and therefore
flat. They give, on the minimal model, arithmetic genus
one for each connected component and Euler number equal
to the number of its rational components. The rational
curve count is at most $b_2$, as stated and used in
\cite[Introduction, p.~283]{DOT2003}.
These are consequences of the proved theorem about
anticanonical divisors; no classification assertion about
class VII surfaces without such a divisor is used.

There is a coefficient issue in the parabolic Inoue case
which the torsion hypothesis resolves explicitly. Write
$E$ for its elliptic curve and $C$ for its reduced rational
cycle. The exact canonical relation is
$K_{Z_0}+E+C=0$; the curves $E$ and $C$ are disjoint,
$E^2=-b_2(Z_0)$, and the intersection matrix of the
components of $C$ has kernel generated by their sum
\cite[Lemma~2.22 and the following paragraph; Proposition~2.27]{Dloussky2006}.
Consequently an effective integral divisor with
$c_1(K_{Z_0}+D_0)_{\mathbb R}=0$ has the form
$D_0=E+kC$, $k\ge0$: intersection with $E$ fixes its
coefficient at one, and the cycle intersection matrix
makes all cycle coefficients equal. Our stronger
torsion relation then gives
$K_{Z_0}+D_0=(k-1)C$ torsion. For any Gauduchon metric
$\omega$, the degree of $\OO(C)$ is
$\int_C\omega>0$, whereas a torsion line has degree
zero. Hence $k=1$. In particular both connected
pieces of the stated reduced support really occur.
This deduction uses the torsion relation, not merely
the vanishing of its real first Chern class.

For a minimal model with $b_2=0$, the nonzero effective
divisor excludes the curve-free Inoue surfaces. The
Hopf alternative has only smooth elliptic curves, which
are pairwise disjoint. This follows as well from the
explicit contraction description used in C.3.8 and its
finite unramified quotient: the preimages of two distinct
curves are disjoint elliptic curves, and the quotient
is locally biholomorphic. All assertions are then
immediate on that model.

We verify the transfer through the actual blowdowns.
Let $\sigma:Z\to Z_0$ be a sequence of point blowups
from a minimal model. A torsion line is represented by
a finite unitary character of the fundamental group,
and a point blowup induces an isomorphism on fundamental
groups. Thus the torsion line $K_Z+D$ descends to $Z_0$.
Pushing forward a meromorphic canonical section with pole
divisor $D$, or its torsion-trivializing power, gives an
effective divisor $D_0$ with $K_{Z_0}+D_0$ torsion.
It is nonzero, since a torsion canonical bundle would
give Kodaira dimension zero. At each individual blowup
the canonical divisor formula and the pushforward section
give the exact divisor equality
\[
 D_{\mathrm{new}}=\sigma^*D_{\mathrm{old}}-E.
\tag{C.VII.2}
\]
The coefficient of $E$ is
$\operatorname{mult}_pD_{\mathrm{old}}-1$.
Effectivity forces the centre to lie on $D_{\mathrm{old}}$.
If this coefficient is positive, the reduced support
is its entire reduced inverse image: a smooth point
acquires a rational leaf, or an ordinary intersection
is subdivided by the exceptional curve. These operations
preserve arithmetic genus and increase both the Euler
number and the rational-component count by one.
If the coefficient is zero, the centre is a smooth point
of a multiplicity-one component and the new support is
its strict transform, isomorphic to the old support.
Thus all the stated support properties transfer in the
required direction. Finally each blowup increases $b_2$
by one and adds precisely its exceptional rational
curve; every other rational curve is the strict transform
of one on the preceding surface. The curve-count
inequality in (C.VII.1) follows.
\end{proof}

\begin{lemma}[C.2.VII.2; Proved global conductor Euler bound]
\label{st:C.2.VII.2}
Let $S$ be a compact irreducible reduced hypersurface in
a smooth complex threefold, with $\omega_S$ torsion and
$\chi(S,\OO_S)=0$. Let $\eta:T\to S$ be normalization,
let $\Gamma$ be its conductor image scheme, and put
$C=\Gamma_{\mathrm{red}}$. Then
\[
 \chi(C,\OO_C)\ge-1.
\tag{C.VII.3}
\]
Suppose in addition that the minimal resolution
$\pi:Z\to T$ is a class VII surface. Then
\[
 e(C)\ge0.
\tag{C.VII.4}
\]
The conductor scheme may be nonreduced and its generic
transverse exponent is unrestricted.
\end{lemma}
\begin{proof}
If the conductor is empty, both conclusions are immediate.
We therefore assume that it is nonempty.
Put $L=\omega_S$ and
$\varepsilon=h^0(S,L)\in\{0,1\}$.
A nonzero section of a torsion line on a compact
connected reduced space trivializes it, by taking a
torsion-trivializing power. Serre duality and
$\chi(\OO_S)=0$ therefore give
$h^1(S,\OO_S)=1+\varepsilon$, with
$\varepsilon=1$ exactly when $L=\OO_S$.
The normalization quotient
$\mathcal Q=\eta_*\OO_T/\OO_S$ has
$H^0(\mathcal Q)\hookrightarrow H^1(S,\OO_S)$,
because $S$ and $T$ have the same constant holomorphic
functions. The actual conductor duality of A.5.12 is
\[
 \mathcal Q=\omega_\Gamma\otimes L^{-1}|_\Gamma .
\]
Serre duality on the Cohen--Macaulay conductor curve
gives $h^1(\Gamma,L|_\Gamma)\le1+\varepsilon$.
The restriction onto $C$ is a quotient of coherent
sheaves on a one-dimensional space. Its kernel has
no second cohomology, so
$h^1(C,L|_C)\le1+\varepsilon$.
Since $L|_C$ is torsion, it has degree zero on every
irreducible component and
$\chi(C,L|_C)=\chi(C,\OO_C)$.
If $\varepsilon=0$, Riemann--Roch gives
$\chi(\OO_C)=h^0(C,L|_C)-h^1(C,L|_C)\ge-1$.
If $\varepsilon=1$, then $L=\OO_S$ and
$h^0(C,\OO_C)\ge1$, while $h^1(C,\OO_C)\le2$.
This proves (C.VII.3), including all nilpotents of
the original conductor scheme in the comparison.

Now assume that $Z$ is class VII. A.5.3 gives a
nonzero effective integral divisor $D$ with
$K_Z+D=(\eta\pi)^*\omega_S$.
Each curve in the reduced strict conductor on $Z$
is rational or elliptic by C.2.VII.1. Its finite image
in $C$ has normalization of genus at most one, by
Riemann--Hurwitz. Let $n$ be the number of irreducible
components of $C$, let $g_i\in\{0,1\}$ be the genera
of their normalizations, and let $r_x$ and $\delta_x$
be the branch number and delta invariant at a singular
point of $C$. The exact normalization sequence and
the topological normalization map give
\[
 \chi(\OO_C)=\sum_i(1-g_i)-\sum_x\delta_x,\qquad
 e(C)=\sum_i(2-2g_i)-\sum_x(r_x-1).
\tag{C.VII.5}
\]
For every reduced curve germ,
$\delta_x\ge r_x-1$: its normalization permits
independent constants on the branches, whereas a
function on the original germ has one common constant.
If $n_0$ is the number of rational-normalization
components, (C.VII.3)--(C.VII.5) imply
$e(C)\ge n_0-1$.

The only possible negative value is consequently
$e(C)=-1$, with every normalized component elliptic.
Then (C.VII.3) says $\sum\delta_x\le1$.
Negativity requires two branches, so there is exactly
one singular point, it has two branches and
$\delta=1$, and the local curve ring is
\[
 \{(a(u),b(v))\in\C\{u\}\oplus\C\{v\}:a(0)=b(0)\}.
\tag{C.VII.6}
\]
Indeed that subring already has colength one in the
normalization, and the curve ring is contained in it
with the same colength. Thus the point is an ordinary
node. Its two branches come either from distinct
elliptic normalizations or from two distinct points
of one elliptic normalization.

We show that this node is incompatible with the
specified hypersurface and resolution maps. At least
one normalization point $t\in T$ over the node receives
preimages of both local branches of $C$.
Otherwise partition the irreducible local branches
of the hypersurface germ $S$ according to the one
conductor branch their normalization receives. Both
parts are nonempty. Two hypersurface branches from
opposite parts intersect in a curve: in the regular
three-dimensional local ambient ring, the ideal of
their two distinct equations has height two.
That intersection curve is contained in the conductor
image and therefore is one of its two branches.
Both surface branches must receive it, contrary to
the partition. With only one local surface branch
the assertion is immediate.

Every $\pi$-exceptional curve not contained in
$D_{\mathrm{red}}$ is disjoint from it. To see this,
write $E$ for such a curve. Minimality of the
resolution and adjunction give $K_Z\cdot E\ge0$,
while $D\cdot E\ge0$, and
$(K_Z+D)\cdot E=0$ forces equality. In particular
$D\cdot E=0$. A resolution fibre over a point of
the conductor must meet its strict transform, by
properness. Since its fibre is connected, this
disjointness proves
$\pi^{-1}(t)\subset D_{\mathrm{red}}$ for the
normalization point just found.

The strict curves over the two elliptic-normalized
branches are elliptic: a rational curve cannot map
nonconstantly onto an elliptic curve. They meet the
connected fibre $\pi^{-1}(t)$. If the curves are
distinct, one connected component of $D_{\mathrm{red}}$
therefore contains two elliptic curves. If they are
the same elliptic curve, their meeting points with
the fibre are distinct. This follows from the exact
factorization of a map from a normal curve through
the normalization of its image: one point cannot
map to both points over the node in (C.VII.6).
The connected resolution fibre consequently gives
a path in $D_{\mathrm{red}}$ joining these distinct
points. Both possibilities contradict C.2.VII.1:
the first has two positive-genus components, and
the second adds a graph cycle to an elliptic
component, giving arithmetic genus at least two.
This excludes the only negative possibility and
proves (C.VII.4).
\end{proof}

\begin{theorem}[C.2.VII.3; Proved sign for class VII normalization]
\label{st:C.2.VII.3}
Let $f:Y\to B$ be a proper connected-fibre holomorphic
map from a smooth compact complex threefold with
$b_2(Y)=0$ to a smooth curve, whose general fibre is a
complex two-torus. Let $mS$ be a full fibre with
irreducible reduced nonnormal support, and suppose
that the minimal resolution $Z$ of its normalization
is class VII. Then
\[
 e(S)\ge0.
\tag{C.VII.7}
\]
No restriction on the singularity type or conductor
multiplicity of $S$ is made.
\end{theorem}
\begin{proof}
The canonical evaluation and adjunction comparison
(A.CD.1) give $\omega_S^{\otimes m}=\OO_S$.
The divisor sequence and Hirzebruch--Riemann--Roch
give $\chi(\OO_S)=0$, because $c_1(\OO_Y(S))=0$
rationally. These arguments require only $b_2(Y)=0$;
the integral condition and Picard-variation assumption
in the later part of A.5.12 are not needed here.
Let $\eta:T\to S$, $\pi:Z\to T$, and
$D>0$ be the maps and divisor of A.5.3. Put
$C=\Gamma_{\mathrm{red}}$ and
$A=(\eta^{-1}C)_{\mathrm{red}}\subset T$.
Thus $K_Z+D$ is torsion, and C.2.VII.1--C.2.VII.2
apply.

We compute $e(T\setminus A)$ through the actual
resolution. As in C.2.VII.2, an exceptional curve
$E$ with coefficient zero in $D$ has
$K_Z\cdot E=D\cdot E=0$.
Adjunction and negative self-intersection then give
$p_a(E)=0$ and $E^2=-2$.
These curves are smooth rational and disjoint
from $D_{\mathrm{red}}$.
Their negative-definite intersection matrix makes
each connected configuration a tree: a repeated
edge or a cycle would give a nonnegative square
for a nonzero vector in the matrix. The intersections
are transverse, since their local intersection
multiplicities are one; three curves cannot meet
at a point, since that would give a triangle.
Connectedness of resolution fibres shows that a
fibre either consists entirely of these curves or
lies entirely in $D_{\mathrm{red}}$.
The zero-coefficient fibres lie outside $A$, since
a fibre over $A$ must meet the strict conductor.

Let $a$ be the total number of zero-coefficient
exceptional curves, and let $b$ be the number of
exceptional fibres lying in $D_{\mathrm{red}}$
whose images in $T$ lie outside $A$. On
$Z\setminus D_{\mathrm{red}}$ the former fibres
are still present and the latter are absent.
A tree with $a_j$ rational curves has Euler
number $a_j+1$, so replacing it by its image
point subtracts $a_j$. Adding back the image
point of an absent positive-coefficient fibre
adds one. Away from these finitely many fibres,
$\pi$ is the required isomorphism of complements.
Additivity of compactly supported Euler
characteristic therefore gives the exact equality
\[
 \begin{aligned}
 e(T\setminus A)
   &=e(Z)-e(D_{\mathrm{red}})-a+b\\
   &=b_2(Z)-r(D_{\mathrm{red}})-a+b
   \ \ge0 .
 \end{aligned}
\tag{C.VII.8}
\]
The inequality follows because the rational curves
counted by $r(D_{\mathrm{red}})$ and by $a$ are
distinct, and their total is at most $b_2(Z)$ by
C.2.VII.1.
Finally $\eta$ is a biholomorphism
$T\setminus A\to S\setminus C$. Hence
\[
 e(S)=e(T\setminus A)+e(C)\ge0
\]
by C.2.VII.2. This proves the assertion for the
original reduced support and all full multiplicities.
\end{proof}

\begin{theorem}[C.2.VII.4; Proved equality comparison]
\label{st:C.2.VII.4}
Retain the irreducible class VII situation of
C.2.VII.3. Let $C=\Gamma_{\rm red}$, let $n_0$
be the number of its rational-normalized components,
and put $v(C)=\sum_x(s_x-1)$, where $s_x$ is the
number of branches at a singular point of $C$.
On the minimal resolution $Z$ let $D$ be the
resolution conductor, let $r$ count the rational
curves in $D_{\rm red}$, and let $a$ count the
zero-coefficient exceptional curves lying over
points outside the conductor. Let $b$ count the
positive-coefficient exceptional fibres over
such points, as in C.2.VII.3. Then $e(S)=0$
is equivalent to the following explicit
normalization and curve data:
\[
 b=0,\qquad b_2(Z)=r+a,
 \qquad
 \begin{cases}
 n_0=0, v(C)=0,\quad\text{or}\\
 n_0=1, v(C)=2.
 \end{cases}
 \tag{C.VII.9}
\]
In the first curve alternative all components
have elliptic normalization and the sum of
their delta invariants is at most one. In the
second the total delta invariant is exactly
two, and every singular point satisfies
$\delta_x=s_x-1$.

If the actual homomorphism
$\Pic^0(Y)\to\Pic^0(Z)$ is nonconstant,
there is the sharper geometric equivalence:
$e(S)=0$ exactly when its normalization is a
smooth minimal Hopf surface and the reduced
conductor image and inverse image are disjoint
unions of smooth elliptic curves, with finite
unramified maps between them.
\end{theorem}
\begin{proof}
The exact comparison already proved in
C.2.VII.3 is
\[
 e(S)=b_2(Z)-r-a+b+e(C),
 \qquad b_2(Z)-r-a\ge0,\quad b\ge0,\quad e(C)\ge0.
 \tag{C.VII.10}
\]
Every normalized component of $C$ has genus
zero or one. Its topological normalization
formula is $e(C)=2n_0-v(C)$. The stronger
bound in C.2.VII.2 is $e(C)\ge n_0-1$.
Thus $e(C)=0$ permits precisely $n_0=0,1$,
with $v(C)=0,2$, respectively. This proves
both directions of (C.VII.9), using actual
branch numbers rather than an assumption
about a local singularity type.
The same proof gives
$\chi(\OO_C)=n_0-\sum\delta_x\ge-1$.
For $n_0=0$ this bounds the total delta
invariant by one. For $n_0=1$, the value
$v(C)=2$ and $\delta_x\ge s_x-1$ force
$\sum\delta_x=2$ and equality at every point.
These are exact consequences for fibres which
satisfy the criterion; no separate claim
that all such abstract curve patterns admit
a torus smoothing is used.

Now suppose that the Picard pullback is
nonconstant. The argument of A.5.12 applies
to $\Pic^0(Y)$ even with only rational
$H^2(Y,\Q)=0$. Namely its nonzero differential
is the actual composite
\[
 H^1(Y,\OO_Y)\longrightarrow H^1(S,\OO_S)
 \longrightarrow H^1(T,\OO_T)
 \longrightarrow H^1(Z,\OO_Z).
 \tag{C.VII.11}
\]
The middle arrow consequently has positive
rank. Put $\varepsilon=h^0(S,\omega_S)$;
the torsion canonical relation makes
$\varepsilon\in\{0,1\}$ and
$h^1(S,\OO_S)=1+\varepsilon$.
The normalization sequence then gives
$h^0(\mathcal Q)\le\varepsilon$.
Conductor duality and restriction onto $C$
therefore give
\[
 h^1(C,\omega_S|_C)\le\varepsilon.
 \tag{C.VII.12}
\]
If $\varepsilon=0$, degree-zero
Riemann--Roch gives $\chi(\OO_C)\ge0$.
If $\varepsilon=1$, the torsion line
$\omega_S$ is trivial, and
$h^0(C,\OO_C)\ge1$ gives the same inequality.
The exact formula (C.INT.26) now implies
that $e(C)=0$ holds precisely when $C$ is
a disjoint union of smooth elliptic curves.

Assume $e(S)=0$. Every rational curve in
$D_{\rm red}$ must be resolution-exceptional,
because its image lies in $C$ and cannot
dominate an elliptic curve. A minimal class
VII model with positive second Betti number
and a torsion anticanonical divisor is one
of the proved configurations used in
C.2.VII.1. An intermediate or
Inoue--Hirzebruch model has no elliptic
curves, and therefore could not dominate
the nonempty elliptic conductor. A
parabolic Inoue model has a rational cycle
of square zero and exactly $b_2$ rational
curves. Its point blowups therefore still
have exactly $b_2(Z)$ rational curves.
The equality $b_2(Z)=r+a$ makes every one
of them resolution-exceptional, including
any zero-coefficient blowup curves. Thus
the full total transform of the minimal
rational cycle, whose square remains zero,
would be exceptional. This contradicts
negative definiteness of a resolution's
exceptional intersection matrix.
Consequently the minimal model is Hopf.
All rational curves on $Z$ are now the
exceptional curves of its point blowups,
and their number is $b_2(Z)$.
The equality $b_2(Z)=r+a$ makes all of
them resolution-exceptional. As in the
last paragraph of C.2.INT.3, a nonempty
blowup sequence would have a last
$(-1)$ curve, contradicting minimality
of the resolution. Thus the normalization
is already a smooth minimal Hopf surface.
Its reduced conductor curves are elliptic;
their finite maps to the smooth elliptic
curves of $C$ are unramified by
Riemann--Hurwitz.
Conversely those normalization and
conductor data have Euler number zero
by the exact normalization formula.
This proves the sharper equivalence.
\end{proof}
